\section{Conclusion}\label{sec:conclusion}

The active-pile chain has an exact arbitrary-state statistic: its expected
number of embedded updates is the sum of the meeting-time means of consecutive
initial interfaces.  The site-clock construction proves scheduler
equivalence, while the compensator converts elapsed interface lifetimes back
to the discrete update count.  A ballot calculation then yields the
double-factorial full-start mean and its $n^{3/2}$ scale.  The finite PGF and
support theorem provide a second description of the hitting law.

The result is path-specific and does not assert independent interfaces.
Generic coalescing-walk and first-passage methods have been subtracted, and
the source search remains bounded.  Accordingly, novelty, priority, posting,
submission, and external release remain on hold.
